\documentclass[../../main.tex]{subfiles}

\begin{document}

\section{归纳定义原理}\label{section:ch01-归纳定义原理}

在讨论归纳定义原理的一般形式之前，首先证明它的一个特殊情形——\cref{lemma:ch01-正整数无限子集}的证明中用到过这种情形。它有助于对一般情形的讨论，使证明中的关键思路更为清晰。

为此，给定 $\mathbb{Z}_{+}$ 的无限子集 $C$，考虑函数 $h: \mathbb{Z}_{+} \longrightarrow C$ 的下述归纳公式：
\begin{align}
h(1) = C \text{ 的最小元},\label{align:ch01-归纳公式特殊}
\end{align}
\begin{align*}
h(i) = [C - h(\{1, \cdots, i-1\})] \text{ 的最小元，对于 } i > 1.
\end{align*}
我们证明存在唯一的一个函数 $h: \mathbb{Z}_{+} \longrightarrow C$ 满足这个归纳公式。

第一步，证明存在定义在 $\mathbb{Z}_{+}$ 的截上的函数满足\eqref{align:ch01-归纳公式特殊}。

\begin{lemma}\label{lemma:ch01-截上归纳}
给定 $n \in \mathbb{Z}_{+}$，存在一个函数
\begin{align*}
f: \{1, \cdots, n\} \longrightarrow C,
\end{align*}
对于定义域中的所有 $i$ 满足\eqref{align:ch01-归纳公式特殊}。
\end{lemma}

\begin{proof}
这个引理讲的是一个依赖于 $n$ 的论断，因此可以用归纳法加以证明。设 $A$ 为所有使引理成立的 $n$ 的集合。我们证明 $A$ 是一个归纳集。从而有 $A = \mathbb{Z}_{+}$。

当 $n = 1$ 时引理成立，因为由
\begin{align*}
f(1) = C \text{ 的最小元}
\end{align*}
定义的函数 $f: \{1\} \longrightarrow C$ 满足\eqref{align:ch01-归纳公式特殊}。

假定引理对于 $n-1$ 成立，以下证明引理对于 $n$ 成立。根据归纳假定，存在函数 $f': \{1, \cdots, n-1\} \longrightarrow C$，对于定义域 $\{1, \cdots, n-1\}$ 中的所有 $i$ 满足\eqref{align:ch01-归纳公式特殊}。用
\begin{align*}
f(i) = f'(i), \quad \text{对于 } i \in \{1, \cdots, n-1\},
\end{align*}
\begin{align*}
f(n) = [C - f'(\{1, \cdots, n-1\})] \text{ 的最小元}
\end{align*}
定义一个函数 $f: \{1, \cdots, n\} \longrightarrow C$。因为 $C$ 为无限集，$f'$ 不是满射；因此集合 $C - f'(\{1, \cdots, n-1\})$ 非空，并且 $f(n)$ 有定义。这个定义是比较容易接受的，它不是用 $f$ 自身而是用给定的函数 $f'$ 来定义 $f$。

容易验证，$f$ 对于定义域中的所有 $i$ 满足\eqref{align:ch01-归纳公式特殊}。因为当 $i \leqslant n-1$ 时，$f$ 等于 $f'$，所以函数 $f'$ 满足\eqref{align:ch01-归纳公式特殊}。当 $i = n$ 时，$f$ 定义为
\begin{align*}
f(n) = [C - f'(\{1, \cdots, n-1\})] \text{ 的最小元},
\end{align*}
而 $f'(\{1, \cdots, n-1\}) = f(\{1, \cdots, n-1\})$，所以 $f$ 也满足\eqref{align:ch01-归纳公式特殊}。
\end{proof}

\begin{lemma}[截上函数的相容性]\label{lemma:ch01-截相容}
设 $f: \{1, \cdots, n\} \longrightarrow C$ 与 $g: \{1, \cdots, m\} \longrightarrow C$ 对于它们各自定义域中的所有 $i$ 都满足\eqref{align:ch01-归纳公式特殊}。则对于两个定义域中所有公共的 $i$，$f(i) = g(i)$。
\end{lemma}

\begin{proof}
设结论不真。令 $i$ 为使得 $f(i) \neq g(i)$ 的最小整数。根据\eqref{align:ch01-归纳公式特殊}
\begin{align*}
f(1) = C \text{ 的最小元} = g(1).
\end{align*}
所以整数 $i$ 不是 $1$。对于所有 $j < i$，我们有 $f(j) = g(j)$。由于 $f$ 和 $g$ 满足\eqref{align:ch01-归纳公式特殊}，所以
\begin{align*}
f(i) = [C - f(\{1, \cdots, i-1\})] \text{ 的最小元},
\end{align*}
\begin{align*}
g(i) = [C - g(\{1, \cdots, i-1\})] \text{ 的最小元}.
\end{align*}
由于 $f(\{1, \cdots, i-1\}) = g(\{1, \cdots, i-1\})$，所以 $f(i) = g(i)$，这与 $i$ 的选取矛盾。
\end{proof}

\begin{theorem}[归纳函数的存在唯一性]\label{theorem:ch01-归纳存在唯一}
存在唯一的一个函数 $h: \mathbb{Z}_{+} \longrightarrow C$，使得\eqref{align:ch01-归纳公式特殊}对于所有 $i \in \mathbb{Z}_{+}$ 成立。
\end{theorem}

\begin{proof}
根据\cref{lemma:ch01-截上归纳}，对于每一个 $n$，存在一个将 $\{1, \cdots, n\}$ 映到 $C$ 中的函数，并且对其定义域中的所有 $i$ 满足\eqref{align:ch01-归纳公式特殊}。给定 $n$，\cref{lemma:ch01-截相容}证明了这样的函数是唯一的，定义域相同的两个这样的函数必定相等。设 $f_{n}: \{1, \cdots, n\} \longrightarrow C$ 表示这个唯一的函数。

现在进行关键的一步。定义一个函数 $h: \mathbb{Z}_{+} \longrightarrow C$，其指派法则是所有 $f_{n}$ 的指派法则的并 $U$。由于 $f_{n}$ 的指派法则是 $\{1, \cdots, n\} \times C$ 的一个子集，因此 $U$ 是 $\mathbb{Z}_{+} \times C$ 的一个子集。我们要证明 $U$ 是函数 $h: \mathbb{Z}_{+} \longrightarrow C$ 的指派法则。

就是说我们要证明 $\mathbb{Z}_{+}$ 的每一个元素 $i$ 恰好是 $U$ 中一个元素的第一个坐标，这是容易的。整数 $i$ 在 $f_{n}$ 定义域中的充分必要条件是 $n \geqslant i$，所以 $U$ 中所有使 $i$ 为其第一个坐标的元素的集合，正好是形如 $(i, f_{n}(i))$ 的所有偶对的集合，其中 $n \geqslant i$。\cref{lemma:ch01-截相容}表明，当 $n, m \geqslant i$ 时 $f_{n}(i) = f_{m}(i)$。因此，$U$ 中所有这些元素都相等，亦即 $U$ 中只有一个元素以 $i$ 为它的第一个坐标。

证明 $h$ 满足条件\eqref{align:ch01-归纳公式特殊}也是容易的，它是以下事实的推论：
\begin{align*}
\text{如果 } i \leqslant n, \text{ 则 } h(i) = f_{n}(i),
\end{align*}
\begin{align*}
f_{n} \text{ 对于定义域中所有 } i \text{ 满足\eqref{align:ch01-归纳公式特殊}.}
\end{align*}
唯一性的证明是\cref{lemma:ch01-截相容}的证明的翻版。
\end{proof}

现在我们正式叙述归纳定义的一般原理。在它的证明中并没有什么新思想，故留作习题。

\begin{theorem}[归纳定义原理]\label{theorem:ch01-归纳定义原理}
设 $A$ 是一个集合，$a_{0}$ 为 $A$ 的一个元素。设 $\rho$ 为一个函数，使得每一个从正整数的一个非空截映到 $A$ 中的函数 $f$ 对应于 $A$ 中一个元素。则存在唯一的一个函数
\begin{align*}
h: \mathbb{Z}_{+} \longrightarrow A,
\end{align*}
使得
\begin{align}
h(1) = a_{0},\label{align:ch01-归纳公式一般}
\end{align}
\begin{align*}
h(i) = \rho(h \mid \{1, \cdots, i-1\}), \quad \text{对于 } i > 1.
\end{align*}
\eqref{align:ch01-归纳公式一般}称为 $h$ 的一个归纳公式（recursion formula）。它决定了 $h(1)$，并且用 $h$ 在所有小于 $i$ 的正整数处的值来表出 $h$ 在 $i > 1$ 处的值。
\end{theorem}

\begin{example}\label{example:ch01-归纳特殊情形}
我们来证明\cref{theorem:ch01-归纳存在唯一}是\cref{theorem:ch01-归纳定义原理}的一个特殊情形。给定 $\mathbb{Z}_{+}$ 的一个无限子集 $C$，设 $a_{0}$ 为 $C$ 的最小元，$\rho$ 定义为：
\begin{align*}
\rho(f) = [C - (f \text{ 的像集})] \text{ 的最小元}.
\end{align*}
由于 $C$ 是一个无限集，$f$ 为从有限集映到 $C$ 中的一个函数，$f$ 的像集不是整个集合 $C$，因而 $\rho$ 的定义是确切的。根据\cref{theorem:ch01-归纳定义原理}，存在一个函数 $h: \mathbb{Z}_{+} \longrightarrow C$，使得 $h(1) = a_{0}$，并且对于 $i > 1$，
\begin{align*}
h(i) = \rho(h \mid \{1, \cdots, i-1\})
\end{align*}
\begin{align*}
= [C - (h \mid \{1, \cdots, i-1\} \text{ 的像集})] \text{ 的最小元}
\end{align*}
\begin{align*}
= [C - h(\{1, \cdots, i-1\})] \text{ 的最小元},
\end{align*}
这正是我们所要证明的。
\end{example}

\begin{example}\label{example:ch01-幂归纳}
给定 $a \in \mathbb{R}$，在第\cref{section:ch01-整数与实数}节的习题中，我们用归纳公式
\begin{align*}
a^{1} = a,
\end{align*}
\begin{align*}
a^{n} = a^{n-1} \cdot a
\end{align*}
“定义”了 $a^{n}$。我们想用\cref{theorem:ch01-归纳定义原理}来严格定义一个函数 $h: \mathbb{Z}_{+} \longrightarrow \mathbb{R}$，使得 $h(n) = a^{n}$。为了应用这个定理，用 $a_{0}$ 表示 $\mathbb{R}$ 的元素 $a$，并且用公式 $\rho(f) = f(m) \cdot a$ 来定义 $\rho$，其中 $f: \{1, \cdots, m\} \longrightarrow \mathbb{R}$。于是存在唯一的一个函数 $h: \mathbb{Z}_{+} \longrightarrow \mathbb{R}$，使得
\begin{align*}
h(1) = a_{0},
\end{align*}
\begin{align*}
h(i) = \rho(h \mid \{1, \cdots, i-1\}), \quad \text{对于 } i > 1.
\end{align*}
这意味着 $h(1) = a$。当 $i > 1$ 时有 $h(i) = h(i-1) \cdot a$。如果以 $a^{i}$ 表示 $h(i)$，则得到了所要证明的
\begin{align*}
a^{1} = a,
\end{align*}
\begin{align*}
a^{i} = a^{i-1} \cdot a.
\end{align*}
\end{example}

\subsection*{习题}

\begin{exercise}\label{exercise:ch01-归纳定义-1}
设 $(b_{1}, b_{2}, \cdots)$ 是实数的一个无穷序列。用归纳法定义它的和 $\sum_{k=1}^{n} b_{k}$ 如下：
\begin{align*}
\sum_{k=1}^{n} b_{k} = b_{1} \quad \text{当 } n = 1,
\end{align*}
\begin{align*}
\sum_{k=1}^{n} b_{k} = \left(\sum_{k=1}^{n-1} b_{k}\right) + b_{n} \quad \text{当 } n > 1.
\end{align*}
设 $A$ 为实数集，选取 $\rho$，使得可以用\cref{theorem:ch01-归纳定义原理}来严格定义这个和。

我们有时用符号 $b_{1} + b_{2} + \cdots + b_{n}$ 表示和 $\sum_{k=1}^{n} b_{k}$。
\end{exercise}

\begin{exercise}\label{exercise:ch01-归纳定义-2}
设 $(b_{1}, b_{2}, \cdots)$ 为实数的一个无穷序列。$\prod_{k=1}^{n} b_{k}$ 的定义为
\begin{align*}
\prod_{k=1}^{1} b_{k} = b_{1},
\end{align*}
\begin{align*}
\prod_{k=1}^{n} b_{k} = \left(\prod_{k=1}^{n-1} b_{k}\right) \cdot b_{n}, \quad \text{当 } n > 1.
\end{align*}
试应用\cref{theorem:ch01-归纳定义原理}来严格定义这个积。我们有时用符号 $b_{1}b_{2}\cdots b_{n}$ 来表示积 $\prod_{k=1}^{n} b_{k}$。
\end{exercise}

\begin{exercise}\label{exercise:ch01-归纳定义-3}
作为习题\cref{exercise:ch01-归纳定义-2}的特例，对于 $n \in \mathbb{Z}_{+}$，给出 $a^{n}$ 与 $n!$ 的定义。
\end{exercise}

\begin{exercise}\label{exercise:ch01-归纳定义-4}
数论中的 Fibonacci 数是用下式归纳定义的：
\begin{align*}
\lambda_{1} = \lambda_{2} = 1,
\end{align*}
\begin{align*}
\lambda_{n} = \lambda_{n-1} + \lambda_{n-2}, \quad \text{对于 } n > 2.
\end{align*}
试用\cref{theorem:ch01-归纳定义原理}给出它的严格定义。
\end{exercise}

\begin{exercise}\label{exercise:ch01-归纳定义-5}
证明存在唯一的一个函数 $h: \mathbb{Z}_{+} \longrightarrow \mathbb{R}_{+}$，满足公式
\begin{align*}
h(1) = 3,
\end{align*}
\begin{align*}
h(i) = [h(i-1) + 1]^{1/2}, \quad \text{对于 } i > 1.
\end{align*}
\end{exercise}

\begin{exercise}\label{exercise:ch01-归纳定义-6}
(a) 证明不存在函数 $h: \mathbb{Z}_{+} \longrightarrow \mathbb{R}_{+}$ 满足公式
\begin{align*}
h(1) = 3,
\end{align*}
\begin{align*}
h(i) = [h(i-1) - 1]^{1/2}, \quad \text{对于 } i > 1.
\end{align*}
试说明为什么这个例子不违背归纳定义原理。

(b) 考虑归纳公式
\begin{align*}
h(1) = 3,
\end{align*}
\begin{align*}
h(i) =
\begin{cases}
[h(i-1) - 1]^{1/2}, & \text{如果 } h(i-1) > 1\\
5, & \text{如果 } h(i-1) \leqslant 1
\end{cases}
\quad \text{对于 } i > 1.
\end{align*}
证明存在唯一的一个函数 $h: \mathbb{Z}_{+} \longrightarrow \mathbb{R}_{+}$ 满足这个公式。
\end{exercise}

\begin{exercise}\label{exercise:ch01-归纳定义-7}
证明\cref{theorem:ch01-归纳定义原理}。
\end{exercise}

\begin{exercise}\label{exercise:ch01-归纳定义-8}
证明归纳定义原理的以下形式：设 $A$ 是一个集合，$\rho$ 是一个函数，使得每一个从正整数 $\mathbb{Z}_{+}$ 的一个截 $S_{n}$ 映到 $A$ 中的函数 $f$ 对应着 $A$ 中一个元素 $\rho(f)$。则存在唯一的一个函数 $h: \mathbb{Z}_{+} \longrightarrow A$，使得对于每一个 $n \in \mathbb{Z}_{+}$，$h(n) = \rho(h \mid S_{n})$。
\end{exercise}

\end{document}